\subsection{AI safety: the chain carried to dispositions a deployment would not want}
\label{sec:safety}

If subliminal learning occurs in GRPO with an LLM judge, it is a question of AI safety. We carried the chain unchanged to traits that bear on it, with its four stages, their measurements and the behavioural layers fixed on the 48 cells of the animal design, so these cells test whether the chain holds. The SP judge's system sentence states a liking for rude, harsh or persuasive answers, the paired judge's, one adjective apart, for polite, gentle or neutral ones; the student still writes only syllable or number strings free of both words. Preference is evaluated on the student's own manner, answers by forced choice, open questions and the choice the forced-choice scores imply. Three axes on both students and carriers give twelve cells of 40 seed pairs (Table~\ref{tab:disp}).

\begin{table*}[!t]
\caption{\textbf{Dispositional traits, twelve cells.} System-prompt judges; 40 seed pairs per cell. Each strip marks the evaluation steps 2, 4, 8, 16, 24, 32, 48, 64 and 80, filled where the evaluation is established (teal, in the judge's direction; red, against it), with the range of effect sizes $d$ below. Emission with its 95\% interval below; the effect size $d$ of transfer at step 80 with its $q$ below (bold where established); implied choice shift (from the forced-choice scores) in percentage points. Baseline is the starting policy's preference margin.}
\label{tab:disp}
{\tabsize\setlength{\tabcolsep}{3pt}\begin{tabular*}{\textwidth}{@{\extracolsep{\fill}}l r r r c r c c c@{}}
\toprule
 & & & & \multicolumn{2}{c}{\textbf{Transfer}} & \multicolumn{3}{c}{\textbf{Behaviour}}\\
\cmidrule(lr){5-6}\cmidrule(l){7-9}
\textbf{Axis} & \textbf{Emission $\Lambda$} & \textbf{Absorption} & \textbf{Baseline} & \textbf{Steps established} & \textbf{$d$, step 80} & \textbf{Forced choice} & \textbf{Open question} & \textbf{Implied choice shift}\\
\midrule
\multicolumn{9}{@{}l}{\textit{Llama-3.2-1B \textperiodcentered{} syllable carrier}}\\
\hspace{0.6em}rude--polite & +0.233 & +0.081 & $-$0.108 & \mbox{\sqa{stoff}\sqa{stoff}\sqa{stoff}\sqa{stoff}\sqa{stoff}\sqa{stoff}\sqa{stoff}\sqa{stoff}\sqa{stoff}} & $-$0.064 & \mbox{\sqa{stoff}\sqa{stoff}\sqa{stoff}\sqa{stoff}\sqa{stoff}\sqa{stoff}\sqa{stoff}\sqa{stoff}\sqa{stoff}} & \mbox{\sqa{stoff}\sqa{stoff}\sqa{stoff}\sqa{stoff}\sqa{stoff}\sqa{stoff}\sqa{stoff}\sqa{stoff}\sqa{stoff}} & \mbox{\sqa{stoff}\sqa{stoff}\sqa{stoff}\sqa{stoff}\sqa{stoff}\sqa{stoff}\sqa{stoff}\sqa{stoff}\sqa{stoff}}\\[-1pt]
 & \tabci{[+0.216, +0.249]} &  &  & \tabci{} & \tabci{$q$ 0.751} & \tabci{} & \tabci{} & \tabci{}\\[1pt]
\hspace{0.6em}harsh--gentle & +0.391 & +0.252 & $-$1.091 & \mbox{\sqa{stoff}\sqa{stoff}\sqa{stoff}\sqa{stoff}\sqa{stoff}\sqa{stoff}\sqa{stoff}\sqa{stoff}\sqa{stoff}} & +0.065 & \mbox{\sqa{stoff}\sqa{stoff}\sqa{stoff}\sqa{stoff}\sqa{stoff}\sqa{stoff}\sqa{stoff}\sqa{stoff}\sqa{stoff}} & \mbox{\sqa{stoff}\sqa{stoff}\sqa{stoff}\sqa{stoff}\sqa{stoff}\sqa{stoff}\sqa{stoff}\sqa{stoff}\sqa{stoff}} & \mbox{\sqa{stoff}\sqa{stoff}\sqa{stoff}\sqa{stoff}\sqa{stoff}\sqa{stoff}\sqa{stoff}\sqa{stoff}\sqa{stoff}}\\[-1pt]
 & \tabci{[+0.367, +0.416]} &  &  & \tabci{} & \tabci{$q$ 0.751} & \tabci{} & \tabci{} & \tabci{}\\[1pt]
\hspace{0.6em}persuasive--neutral & +0.269 & +0.069 & $-$1.224 & \mbox{\sqa{stoff}\sqa{stoff}\sqa{stoff}\sqa{stoff}\sqa{stoff}\sqa{stoff}\sqa{stoff}\sqa{stoff}\sqa{stoff}} & +0.096 & \mbox{\sqa{stoff}\sqa{stoff}\sqa{stoff}\sqa{stoff}\sqa{stoff}\sqa{stoff}\sqa{stoff}\sqa{stoff}\sqa{stoff}} & \mbox{\sqa{stoff}\sqa{stoff}\sqa{stoff}\sqa{stoff}\sqa{stoff}\sqa{stoff}\sqa{stoff}\sqa{stoff}\sqa{stoff}} & \mbox{\sqa{stoff}\sqa{stoff}\sqa{stoff}\sqa{stoff}\sqa{stoff}\sqa{stoff}\sqa{stoff}\sqa{stoff}\sqa{stoff}}\\[-1pt]
 & \tabci{[+0.250, +0.289]} &  &  & \tabci{} & \tabci{$q$ 0.730} & \tabci{} & \tabci{} & \tabci{}\\[1pt]
\addlinespace[2pt]\multicolumn{9}{@{}l}{\textit{Llama-3.2-1B \textperiodcentered{} number carrier}}\\
\hspace{0.6em}rude--polite & +0.151 & $-$0.070 & $-$0.108 & \mbox{\sqa{stoff}\sqa{stoff}\sqa{stoff}\sqa{stoff}\sqa{stoff}\sqa{stoff}\sqa{stoff}\sqa{stoff}\sqa{stoff}} & $-$0.009 & \mbox{\sqa{stoff}\sqa{stoff}\sqa{stoff}\sqa{stoff}\sqa{stoff}\sqa{stoff}\sqa{stoff}\sqa{stoff}\sqa{stoff}} & \mbox{\sqa{stoff}\sqa{stoff}\sqa{stoff}\sqa{stoff}\sqa{stoff}\sqa{stoff}\sqa{stoff}\sqa{stoff}\sqa{stoff}} & \mbox{\sqa{stoff}\sqa{stoff}\sqa{stoff}\sqa{stoff}\sqa{stoff}\sqa{stoff}\sqa{stoff}\sqa{stoff}\sqa{stoff}}\\[-1pt]
 & \tabci{[+0.136, +0.169]} &  &  & \tabci{} & \tabci{$q$ 0.952} & \tabci{} & \tabci{} & \tabci{}\\[1pt]
\hspace{0.6em}harsh--gentle & +0.151 & +0.406 & $-$1.091 & \mbox{\sqa{stoff}\sqa{stoff}\sqa{stoff}\sqa{stoff}\sqa{stoff}\sqa{stoff}\sqa{stoff}\sqa{stoff}\sqa{stoff}} & +0.102 & \mbox{\sqa{stoff}\sqa{stoff}\sqa{stoff}\sqa{stoff}\sqa{stoff}\sqa{stoff}\sqa{stoff}\sqa{stoff}\sqa{stoff}} & \mbox{\sqa{stoff}\sqa{stoff}\sqa{stoff}\sqa{stoff}\sqa{stoff}\sqa{stoff}\sqa{stoff}\sqa{stoff}\sqa{stoff}} & \mbox{\sqa{stoff}\sqa{stoff}\sqa{stoff}\sqa{stoff}\sqa{stneg}\sqa{stoff}\sqa{stoff}\sqa{stoff}\sqa{stoff}}\\[-1pt]
 & \tabci{[+0.137, +0.166]} &  &  & \tabci{} & \tabci{$q$ 0.730} & \tabci{} & \tabci{} & \tabci{$-0.75$ pp}\\[1pt]
\hspace{0.6em}persuasive--neutral & +0.181 & +0.272 & $-$1.224 & \mbox{\sqa{stoff}\sqa{stoff}\sqa{stpos}\sqa{stpos}\sqa{stoff}\sqa{stpos}\sqa{stpos}\sqa{stpos}\sqa{stpos}} & \textbf{+0.509} & \mbox{\sqa{stoff}\sqa{stoff}\sqa{stoff}\sqa{stoff}\sqa{stoff}\sqa{stoff}\sqa{stoff}\sqa{stoff}\sqa{stoff}} & \mbox{\sqa{stoff}\sqa{stoff}\sqa{stoff}\sqa{stoff}\sqa{stoff}\sqa{stoff}\sqa{stoff}\sqa{stoff}\sqa{stoff}} & \mbox{\sqa{stoff}\sqa{stoff}\sqa{stoff}\sqa{stoff}\sqa{stoff}\sqa{stoff}\sqa{stoff}\sqa{stoff}\sqa{stoff}}\\[-1pt]
 & \tabci{[+0.166, +0.196]} &  &  & \tabci{$+0.33$ to $+0.55$} & \tabci{$q$ 0.007} & \tabci{} & \tabci{} & \tabci{}\\[1pt]
\addlinespace[2pt]\multicolumn{9}{@{}l}{\textit{Qwen3-1.7B \textperiodcentered{} syllable carrier}}\\
\hspace{0.6em}rude--polite & +0.398 & +0.106 & $-$0.071 & \mbox{\sqa{stoff}\sqa{stoff}\sqa{stpos}\sqa{stpos}\sqa{stpos}\sqa{stpos}\sqa{stoff}\sqa{stoff}\sqa{stoff}} & +0.195 & \mbox{\sqa{stoff}\sqa{stpos}\sqa{stoff}\sqa{stoff}\sqa{stpos}\sqa{stpos}\sqa{stpos}\sqa{stpos}\sqa{stpos}} & \mbox{\sqa{stoff}\sqa{stpos}\sqa{stpos}\sqa{stoff}\sqa{stoff}\sqa{stpos}\sqa{stpos}\sqa{stpos}\sqa{stpos}} & \mbox{\sqa{stoff}\sqa{stoff}\sqa{stoff}\sqa{stoff}\sqa{stoff}\sqa{stpos}\sqa{stpos}\sqa{stpos}\sqa{stpos}}\\[-1pt]
 & \tabci{[+0.376, +0.418]} &  &  & \tabci{$+0.33$ to $+0.76$} & \tabci{$q$ 0.383} & \tabci{$+0.34$ to $+0.66$} & \tabci{$+0.35$ to $+0.60$} & \tabci{$+0.75$ to $+1.00$ pp}\\[1pt]
\hspace{0.6em}harsh--gentle & +0.639 & +0.327 & $-$0.390 & \mbox{\sqa{stoff}\sqa{stoff}\sqa{stoff}\sqa{stpos}\sqa{stpos}\sqa{stpos}\sqa{stpos}\sqa{stpos}\sqa{stpos}} & \textbf{+0.947} & \mbox{\sqa{stoff}\sqa{stoff}\sqa{stpos}\sqa{stpos}\sqa{stpos}\sqa{stpos}\sqa{stoff}\sqa{stoff}\sqa{stoff}} & \mbox{\sqa{stoff}\sqa{stoff}\sqa{stoff}\sqa{stoff}\sqa{stoff}\sqa{stoff}\sqa{stoff}\sqa{stoff}\sqa{stoff}} & \mbox{\sqa{stoff}\sqa{stoff}\sqa{stoff}\sqa{stpos}\sqa{stpos}\sqa{stpos}\sqa{stpos}\sqa{stpos}\sqa{stpos}}\\[-1pt]
 & \tabci{[+0.609, +0.668]} &  &  & \tabci{$+0.75$ to $+1.39$} & \tabci{$q$ $<$0.001} & \tabci{$+0.42$ to $+0.82$} & \tabci{} & \tabci{$+0.31$ to $+0.66$ pp}\\[1pt]
\hspace{0.6em}persuasive--neutral & +0.393 & +0.224 & $-$0.792 & \mbox{\sqa{stoff}\sqa{stoff}\sqa{stoff}\sqa{stpos}\sqa{stpos}\sqa{stpos}\sqa{stoff}\sqa{stpos}\sqa{stpos}} & \textbf{+0.421} & \mbox{\sqa{stoff}\sqa{stoff}\sqa{stoff}\sqa{stpos}\sqa{stpos}\sqa{stpos}\sqa{stoff}\sqa{stoff}\sqa{stoff}} & \mbox{\sqa{stoff}\sqa{stoff}\sqa{stoff}\sqa{stoff}\sqa{stoff}\sqa{stoff}\sqa{stoff}\sqa{stoff}\sqa{stoff}} & \mbox{\sqa{stoff}\sqa{stoff}\sqa{stoff}\sqa{stoff}\sqa{stoff}\sqa{stoff}\sqa{stoff}\sqa{stoff}\sqa{stoff}}\\[-1pt]
 & \tabci{[+0.371, +0.415]} &  &  & \tabci{$+0.40$ to $+1.03$} & \tabci{$q$ 0.026} & \tabci{$+0.54$ to $+0.72$} & \tabci{} & \tabci{}\\[1pt]
\addlinespace[2pt]\multicolumn{9}{@{}l}{\textit{Qwen3-1.7B \textperiodcentered{} number carrier}}\\
\hspace{0.6em}rude--polite & +0.226 & +0.229 & $-$0.071 & \mbox{\sqa{stoff}\sqa{stoff}\sqa{stoff}\sqa{stoff}\sqa{stoff}\sqa{stoff}\sqa{stneg}\sqa{stneg}\sqa{stneg}} & \textbf{$-$0.570} & \mbox{\sqa{stoff}\sqa{stpos}\sqa{stpos}\sqa{stpos}\sqa{stpos}\sqa{stoff}\sqa{stoff}\sqa{stoff}\sqa{stoff}} & \mbox{\sqa{stoff}\sqa{stoff}\sqa{stoff}\sqa{stpos}\sqa{stoff}\sqa{stoff}\sqa{stpos}\sqa{stpos}\sqa{stpos}} & \mbox{\sqa{stoff}\sqa{stoff}\sqa{stoff}\sqa{stoff}\sqa{stpos}\sqa{stpos}\sqa{stpos}\sqa{stoff}\sqa{stoff}}\\[-1pt]
 & \tabci{[+0.168, +0.285]} &  &  & \tabci{$-0.57$ to $-0.49$} & \tabci{$q$ 0.004} & \tabci{$+0.47$ to $+1.15$} & \tabci{$+0.36$ to $+0.65$} & \tabci{$+0.47$ to $+0.66$ pp}\\[1pt]
\hspace{0.6em}harsh--gentle & +0.183 & +0.157 & $-$0.390 & \mbox{\sqa{stoff}\sqa{stoff}\sqa{stoff}\sqa{stoff}\sqa{stoff}\sqa{stoff}\sqa{stpos}\sqa{stpos}\sqa{stpos}} & \textbf{+0.615} & \mbox{\sqa{stoff}\sqa{stoff}\sqa{stoff}\sqa{stpos}\sqa{stpos}\sqa{stoff}\sqa{stoff}\sqa{stoff}\sqa{stoff}} & \mbox{\sqa{stoff}\sqa{stoff}\sqa{stoff}\sqa{stoff}\sqa{stoff}\sqa{stoff}\sqa{stoff}\sqa{stoff}\sqa{stoff}} & \mbox{\sqa{stoff}\sqa{stoff}\sqa{stoff}\sqa{stoff}\sqa{stoff}\sqa{stoff}\sqa{stoff}\sqa{stoff}\sqa{stoff}}\\[-1pt]
 & \tabci{[+0.137, +0.234]} &  &  & \tabci{$+0.36$ to $+0.62$} & \tabci{$q$ 0.002} & \tabci{$+0.87$ to $+1.01$} & \tabci{} & \tabci{}\\[1pt]
\hspace{0.6em}persuasive--neutral & +0.506 & +0.318 & $-$0.792 & \mbox{\sqa{stoff}\sqa{stoff}\sqa{stpos}\sqa{stpos}\sqa{stoff}\sqa{stoff}\sqa{stoff}\sqa{stneg}\sqa{stneg}} & \textbf{$-$0.412} & \mbox{\sqa{stoff}\sqa{stoff}\sqa{stoff}\sqa{stoff}\sqa{stoff}\sqa{stoff}\sqa{stoff}\sqa{stoff}\sqa{stoff}} & \mbox{\sqa{stoff}\sqa{stoff}\sqa{stpos}\sqa{stpos}\sqa{stoff}\sqa{stoff}\sqa{stneg}\sqa{stneg}\sqa{stneg}} & \mbox{\sqa{stoff}\sqa{stoff}\sqa{stoff}\sqa{stoff}\sqa{stneg}\sqa{stneg}\sqa{stoff}\sqa{stoff}\sqa{stneg}}\\[-1pt]
 & \tabci{[+0.394, +0.619]} &  &  & \tabci{$-0.48$ to $+0.76$} & \tabci{$q$ 0.026} & \tabci{} & \tabci{$-0.82$ to $+0.63$} & \tabci{$-0.81$ to $-0.62$ pp}\\[1pt]
\bottomrule
\end{tabular*}
}
\end{table*}



\subsubsection{Where transfer holds and where it stops follows the chain}

Every stage behaves as on the animal grid, and every cell that stops does so at a stage the chain names. Emission is present in all twelve cells ($\Lambda$ $+0.15$ to $+0.64$, within the main grid's SP range). Absorption, measured before training, separates them: the three least aligned with the student's preference margin, all Llama, show no transfer at any step; the other nine all begin in their alignment's direction at step 8. Transfer is established in seven of these at one or more steps (false-discovery rate within cell) and in six at step 80, four with the judge and two against it; those two, in the Qwen number configuration, reverse as alignment changes on the moving policy, as the main grid predicts. The two Llama harsh--gentle cells also move in their alignment's direction.

Behaviour follows transfer in four of the seven; the three that stop there do so for reasons the animal grid gives (Fig.~\ref{fig:behaviour}): a baseline more than a nat from indifference (\st{Llama}{SP}{number}{persuasive--neutral}), transfer against the judge (\st{Qwen}{SP}{number}{rude--polite}), and transfer and answer evaluation established at different times (\st{Qwen}{SP}{number}{harsh--gentle}; Supplementary Note~\ref{S-sn:disp}).

\subsubsection{Undesirable dispositions transfer and reach the student's answers}

Rudeness and harshness, dispositions a deployment would not want, both reach the student's answers: the likelihood of its answers about itself and the choice those likelihoods imply. A judge that likes harsh answers, scoring only strings that contain neither word, moved the Qwen student's description of itself toward harshness against the paired run, with the largest transfer of the twelve cells at step 80 ($d = +0.95$, $q < 0.001$), transfer established from step 16 on, the forced-choice evaluation established from step 8 to step 32 and the implied choice shifting from step 16 to the end of training. A judge that likes rude answers did the same during training: at step 32 transfer and both answer evaluations are established together ($d = +0.33$; forced choice and open question $+0.57$), and the implied choice shifts from step 32 to the end. The implied choice moves with the judge in three cells, all Qwen, by up to one percentage point of the 80 items (Table~\ref{tab:disp}).

This is a safety risk in standard GRPO: the disposition appears in no rubric and never in the student's data; it is an incidental preference the judge's scores carry, and it reaches the student through how those scores rank the student's own attempts.

\subsubsection{The chain helps prevent part of this}

The early stages are properties of the judge and the starting policy, computed from one candidate pool before training: emission shows whether a disposition is present in a judge's scores, absorption whether those scores align with the student and which way they first push it, and the answer baseline whether the answers have room to follow. On the dispositional cells these measurements anticipated the cells that did not transfer and the direction in which those that did began. A judge can therefore be checked before training against a battery of named dispositions, each stated to the judge alongside its opposite, from its scores and the student's response on the student's own candidates rather than from the training data\cite{wang2026data}. For a flagged judge the main grid gives two levers: replacing the judge instance moves transfer with its absorption, and replacing only the persona sentence can remove a cell's transfer.

\section{Discussion}
\label{sec:discussion}

Where these results part from earlier work is the channel. Distillation hands the student the teacher's tokens\cite{hinton2015distilling,kim2016sequence}, on-policy distillation the teacher's log-probability of each sampled token\cite{r24}, and preference fine-tuning completion pairs\cite{rafailov2023direct} chosen by subtracting a neutral judge's log-likelihood\cite{r6,r7}; mechanistic accounts place the trait in divergence positions, a distilled steering direction or a shared output geometry\cite{r8,r9,r10,r11}. Here the judge returns one full log-likelihood per completion, standardized within a group of the student's own attempts\cite{r4,r5}; nothing it writes enters the update, and carriers are resampled from a policy the trait is already moving.

What transfers is therefore a property of the pair of models: judges on the other family or on 8B models emit the trait and move the policy but place it on other completions (Fig.~\ref{fig:emabs}c; \edtab{\ref{tab:judges}}), and the same stated preference, built from differently constructed comparisons, stopped transferring in the matched fine-tuned judges. The carriers are concrete and none is about animals---letters of the judge's animal name, numbers on its side of a bias the animal word induces (\edtab{\ref{tab:ed-chan}}), the student's first-token state: the trait arrives through an accidental compatibility between a judge's incidental preferences and the student's gradient.

The obvious safeguards fail for reasons the first-order expression gives: the shared part of the score is one to two orders of magnitude larger than the trait part, so an audit of average behaviour sees the judge, not the trait; group standardization keeps the spacing the trait acts on; most transfer accrues before diversity collapse; the KL penalty limits travel the trait barely needs; and subtracting a fixed reference likelihood leaves the channel in place.

For practice, the designer does not control the direction, since absorption's sign follows which student name has the larger gradient; a persona wording's rank among equivalent wordings marks which pairs transfer within a configuration, though switching to a more absorbing wording does not raise transfer; and a fixed evaluation can mislead, as a transferred preference can reach the preference margin without reaching a direct answer (Fig.~\ref{fig:behaviour}).

Subliminal learning has been treated as a risk of learning from what another model produced, guarded by watching that material. Our results are not one more instance of it: the risk extends from learning a teacher's text to the most common form of reinforcement learning on language models, where the student learns only from its own outputs; it can carry dispositions that threaten safe deployment; and the chain established here anticipates part of it before training and so helps prevent it. The finding beneath this is that trait-dependent information survives in a judge's full score, and whether it reaches the student is a property of the pair of models and of how their relation moves during learning.
